\documentclass[a4paper,11pt]{article}
\usepackage[utf8]{inputenc}
\usepackage[T1]{fontenc}
\usepackage[ngerman]{babel}
\usepackage{lmodern}\usepackage{microtype}
\usepackage{amsmath,amssymb,bm,mathtools}
\usepackage{geometry}
\geometry{a4paper,left=2.3cm,right=2.3cm,top=2.0cm,bottom=2.0cm}
\usepackage{booktabs,array}
\usepackage{xcolor}\definecolor{bokug}{RGB}{0,135,60}
\usepackage{tikz}
\usetikzlibrary{patterns,arrows.meta,calc,decorations.pathmorphing}
\usepackage{fancyhdr}\pagestyle{fancy}
\renewcommand{\headrulewidth}{0.4pt}
\usepackage{enumitem}\setlength{\parindent}{0pt}\setlength{\parskip}{4pt}

\begin{document}
\fancyhead[L]{\small VU Mechanik – LAWI 100515 (PI)}
\fancyhead[R]{\small SS 2026}
\fancyfoot[C]{\thepage}
\begin{center}{\large\bfseries\color{bokug}Universität für Bodenkultur Wien}\\[2pt]{\large\bfseries VU Mechanik – LAWI 100515 (PI)}\\[4pt]Pr\"ufung \textbf{Gruppe A} \hfill Datum: 26.06.2026\\[2pt]\rule{\linewidth}{0.4pt}\end{center}

\textbf{Gegebenes System}\\[-2pt]
Das System besteht aus zwei Scheiben, die durch ein Gelenk in $C$ verbunden sind: einem Biegetr\"ager $A$--$C$ (feste Einspannung in $A$) und einem Fachwerk (Rollenlager in $B$).
\begin{center}
\begin{tikzpicture}[scale=1.15, rotate=90, >=Stealth, line join=round]
  \coordinate (A) at (0.000,0.000);
  \coordinate (C) at (4.000,0.000);
  \coordinate (O1) at (6.000,0.000);
  \coordinate (O2) at (8.000,0.000);
  \coordinate (U0) at (5.000,-2.000);
  \coordinate (U1) at (7.000,-2.000);
  \draw[thick] (C) -- (O1);
  \draw[thick] (O1) -- (O2);
  \draw[thick] (U0) -- (U1);
  \draw[thick] (C) -- (U0);
  \draw[thick] (U0) -- (O1);
  \draw[thick] (O1) -- (U1);
  \draw[thick] (U1) -- (O2);
  \draw[line width=2.2pt] (A) -- (C);
  \fill (C) circle (1.6pt);
  \fill (O1) circle (1.6pt);
  \fill (O2) circle (1.6pt);
  \fill (U0) circle (1.6pt);
  \fill (U1) circle (1.6pt);
  \draw[fill=white,line width=1pt] (C) circle (3.2pt);
  \draw[line width=1.1pt] (0.000,-0.550) -- (0.000,0.550);
  \fill[pattern=north east lines] (-0.220,-0.550) rectangle (0.000,0.550);
  \draw[line width=1.1pt] (-0.220,-0.550) -- (-0.220,0.550);
  \draw[thick] (8.000,0.000) -- (7.780,-0.340) -- (8.220,-0.340) -- cycle;
  \draw[thick] (7.840,-0.460) circle (0.06) (8.160,-0.460) circle (0.06);
  \draw[line width=1.1pt] (7.700,-0.540) -- (8.300,-0.540);
  \fill[pattern=north east lines] (7.700,-0.700) rectangle (8.300,-0.540);
  \draw[blue!70!black,thick] (0.000,0.550) -- (4.000,0.550);
  \draw[->,blue!70!black] (0.000,0.550) -- (0.000,0.06);
  \draw[->,blue!70!black] (0.667,0.550) -- (0.667,0.06);
  \draw[->,blue!70!black] (1.333,0.550) -- (1.333,0.06);
  \draw[->,blue!70!black] (2.000,0.550) -- (2.000,0.06);
  \draw[->,blue!70!black] (2.667,0.550) -- (2.667,0.06);
  \draw[->,blue!70!black] (3.333,0.550) -- (3.333,0.06);
  \draw[->,blue!70!black] (4.000,0.550) -- (4.000,0.06);
  \node[blue!70!black,above] at (2.000,0.550) {$q$};
  \draw[->,red!80!black,line width=1.1pt] (7.000,-1.050) -- (7.000,-2.000);
  \node[red!80!black,above right] at (7.000,-1.450) {$P$};
  \draw[->,red!80!black,line width=1.1pt] (6.050,-2.000) -- (7.000,-2.000);
  \node[red!80!black,below] at (6.450,-2.000) {$H$};
  \node[circle,fill=white,draw=gray!60,inner sep=0.7pt,font=\scriptsize] at (5.000,0.000) {\textbf{1}};
  \node[circle,fill=white,draw=gray!60,inner sep=0.7pt,font=\scriptsize] at (7.000,0.000) {\textbf{2}};
  \node[circle,fill=white,draw=gray!60,inner sep=0.7pt,font=\scriptsize] at (6.000,-2.000) {\textbf{3}};
  \node[circle,fill=white,draw=gray!60,inner sep=0.7pt,font=\scriptsize] at (4.500,-1.000) {\textbf{4}};
  \node[circle,fill=white,draw=gray!60,inner sep=0.7pt,font=\scriptsize] at (5.500,-1.000) {\textbf{5}};
  \node[circle,fill=white,draw=gray!60,inner sep=0.7pt,font=\scriptsize] at (6.500,-1.000) {\textbf{6}};
  \node[circle,fill=white,draw=gray!60,inner sep=0.7pt,font=\scriptsize] at (7.500,-1.000) {\textbf{7}};
  \node[above left=1pt] at (A) {$A$};
  \node[above=4pt] at (C) {$C$};
  \node[above=2pt] at (O1) {$O_{1}$};
  \node[above right=2pt] at (O2) {$B$};
  \node[below=2pt] at (U0) {$U_{0}$};
  \node[below=2pt] at (U1) {$U_{1}$};
  \draw[<->,gray] (0.000,-2.950) -- (4.000,-2.950) node[midway,below,font=\footnotesize]{$a=4$\,m};
  \draw[<->,gray] (4.000,-2.950) -- (6.000,-2.950) node[midway,below,font=\footnotesize]{$\ell=2$\,m};
  \draw[<->,gray] (3.550,0) -- (3.550,-2.000) node[midway,left,font=\footnotesize]{$h=2$\,m};
\end{tikzpicture}
\end{center}
\textbf{Gegebene Gr\"o\ss en:}\quad $a = 4$\,m\quad $\ell = 2$\,m\quad $h = 2$\,m\quad $q = 3$\,kN/m\quad $P = 12$\,kN\quad $H = 8$\,kN
\par\medskip
\textbf{Aufgabenstellung}\par
\begin{enumerate}[label=\textbf{\arabic*.}]
  \item Bestimmen Sie den Grad der statischen Bestimmtheit des Systems.
  \item Berechnen Sie s\"amtliche Auflagerreaktionen sowie die Gelenkkraft in $C$.
  \item Ermitteln Sie alle Stabkr\"afte des Fachwerks (Zug/Druck angeben).
  \item Zeichnen Sie die Verl\"aufe der Schnittgr\"o\ss en $N(x)$, $V(x)$ und $M(x)$ f\"ur den Biegetr\"ager $A$--$C$.
\end{enumerate}
\end{document}